\Article{14}
Having described the relations of the horizontal magnetic force geometrically, we now develop a way to calculate them. On the Earth's surface $V$ becomes a function of just two variables. We take these to be geographical longitude reckoned eastward from an arbitrary initial meridian and the distance from the north pole. We denote the former by $\lambda$ and the latter, the complement of geographical latitude, by $u$. If we regard the Earth as formed by rotating an ellipse with semimajor axis $=R$ and semiminor axis $=(1-\varepsilon)R$ about the latter axis, the length of an element of the meridian is
\SourceLeaf{142}
\[
=\frac{(1-\varepsilon)^2R\cdot\mathrm{d}u}{(1-(2\varepsilon-\varepsilon\varepsilon)\cos u^2)^{\frac32}}
\]
and the length of an element of the parallel is
\[
=\frac{R\sin u\cdot\mathrm{d}\lambda}{\BareRoot{(1-(2\varepsilon-\varepsilon\varepsilon)\cos u^2)}}
\]

If the horizontal magnetic force is resolved into two parts, one, $X$, acting along the terrestrial meridian and the other, $Y$, perpendicular to it, and if $X$ is taken as positive when it points north and $Y$ as positive when it points west, we obtain
\[
\begin{aligned}
X&=-\frac{(1-(2\varepsilon-\varepsilon\varepsilon)\cos u^2)^{\frac32}}{(1-\varepsilon)^2}\cdot\frac{\mathrm{d}V}{R\mathrm{d}u}\\
Y&=-\BareRoot{(1-(2\varepsilon-\varepsilon\varepsilon)\cos u^2)}\cdot\frac{\mathrm{d}V}{R\sin u\cdot\mathrm{d}\lambda}
\end{aligned}
\]
The total horizontal force is then
\[
=\BareRoot{(XX+YY)}
\]
and the tangent of the declination is
\[
=\frac YX
\]

If the square of the flattening $\varepsilon$ is neglected, these expressions become
\[
\begin{aligned}
X&=-(1+(2-3\cos u^2)\varepsilon)\cdot\frac{\mathrm{d}V}{R\mathrm{d}u}\\
Y&=-(1-\varepsilon\cos u^2)\cdot\frac{\mathrm{d}V}{R\sin u\cdot\mathrm{d}\lambda}
\end{aligned}
\]
or, if the flattening is left out altogether,
\[
X=-\frac{\mathrm{d}V}{R\mathrm{d}u}\qquad Y=-\frac{\mathrm{d}V}{R\sin u\cdot\mathrm{d}\lambda}
\]

The observations available so far are still far too scanty, and most of them far too rough, to make it advisable at present to account for the spheroidal shape of the Earth. Doing so would pose no difficulty in itself, but it would \SourceLeaf{143}greatly complicate the calculations to no benefit. We shall therefore keep to the formulas last given and treat the Earth as a sphere whose radius is $=R$.

